\section{Architecture}

\subsection*{System Architecture Overview}

The system follows a client-server architecture:

\begin{itemize}
    \item The frontend is implemented using React and TypeScript
    \item The backend is implemented using Kotlin and Spring
    \item Communication occurs through REST APIs
\end{itemize}


\subsection{How Much Up-Front?} %TODO endre tittel?  
Based on our experience in the project, we performed relatively little upfront architectural work.
This was mainly due to time constraints, learning new technologies, and a focus on delivering working functionality early in the development
This aligns with agile principles, where teams often aim to design “just enough” architecture to get started.

However, when reflecting on the paper\cite{UpFront}, our approach may have been leaning too far towards minimal upfront design.
The paper highlights that too little upfront planning can lead to an “accidental architecture”,
where decisions are made reactively rather than strategically.
While we have not yet encountered major issues caused by this,
there are clear areas where this could become problematic in later stages of the project.

For example, decisions regarding authentication and integration with external systems, such as using an existing student login solution,
have not been planned in advance.
Introducing such functionality later may require significant changes to both frontend and backend structure, indicating a potential architectural gap.

At the same time, our project included some factors such as needing to deliver functionality early, which supports a more emergent architecture approach.
According to the paper, this justifies reducing upfront work in order to remain flexible and responsive to change. 

In hindsight, a slightly higher level of upfront architectural planning would likely have been beneficial,
especially in identifying external requirements and defining key system boundaries earlier.
This could have reduced the risk of future refactoring while still maintaining flexibility.

Overall, our architectural decisions were discussed, but not always systematically planned.
The project reflects the trade-off described in the paper.
Although our current approach has worked so far, a more balanced approach between upfront planning and emergent design would likely improve the system as it expands.


\subsection{Architectural Concerns}
In our project we don’t have a defined software architect.
Instead we share the responsibility across the team, where everyone is expected to reflect on whether what they implement fits the system.
The paper “\textit{What do software architects really do?}”\cite{SoftwareArchitectsReallyDo} suggests that architectural work includes both an “inwards” and “outwards” focus,
and that there should be a balance between them.

After the first assignment we did interviews and created personas, which helped us understand user needs.
This matches the inwards focus from the paper.
We then had a meeting to plan a solution and started building a basic frontend and backend.
We also spent time defining data structures and setting up an API, which shows some focus on architecture design.

The paper suggests a 50:25:25 split between internal design work, getting input, and communicating the architecture outwards.
However, we likely didn’t have an ideal balance, especially in the short term.

One of the biggest challenges we experienced is what the paper describes as “firefighting”,
where immediate problems take priority over larger architectural concerns.
Some decisions (like simple data structures and API design) were made quickly to get a working prototype.
While this helped progress, it can introduce technical debt and make future scaling or changes harder.

Looking at the project over time, our focus has shifted.
In the beginning, most of the work was external,
since we needed input to understand the problem.
In later sprints, the focus has been more internal,
since we needed to implement functionality.
However, this also meant less continuous feedback from stakeholders,
which could lead to misalignment with user needs.

Overall, we did consider architectural concerns, but not in a fully balanced way.
We might be slightly leaning towards the absent architect antipattern, since architectural responsibility is
shared and sometimes overshadowed by implementation work or firefighting.

It should also be noted that the 50:25:25 split is based on experience and not necessarily a strict rule.
For a team like ours, with less experience and shared responsibilities, a different balance may be more ideal.
